\documentclass{article}
\usepackage{lmodern}
\usepackage{amsmath}
\usepackage{listings}
\usepackage{fullpage}
\usepackage{parskip}
\usepackage{xcolor}
\lstset{
	basicstyle=\ttfamily,
	columns=fullflexible,
	frame=single,
	breaklines=true,
	postbreak=\mbox{\textcolor{red}{$\hookrightarrow$}\space},
}
\begin{document}
\title{Experiment `p\_4\_2\_c\_cumulative\_totals\_seq' Results}
\author{\today}
\date{}
\maketitle
\textbf{Experiment outcome: }SUCCESS
\\\textbf{Bad responses: }0
\\\textbf{Responses containing}~\texttt{assume}~\textbf{: }0
\\\textbf{Responses containing}~\texttt{expect}~\textbf{: }0
\\\textbf{Resolution attempts: }2
\\\textbf{Hard fails (resolution): }0
\\\textbf{Soft fails (resolution): }0
\\\textbf{Verification attempts: }2
\section*{Problem Specification}
\textbf{Problem name: }p\_4\_2\_c\_cumulative\_totals\_seq
\\\textbf{Natural language statement: }null
\\\textbf{Method signature: }p\_4\_2\_c\_cumulative\_totals\_seq(inputs: seq<int>) returns (totals: seq<int>)
\subsection*{Ensures}
\begin{itemize}
\item \texttt{|totals| == |inputs|}
\item \texttt{forall i :: 0 <= i < |inputs| ==> totals[i] == (if i == 0 then inputs[0] else totals[i - 1] + inputs[i])}
\end{itemize}
\subsection*{Requires}
\begin{itemize}
\item \texttt{0 < |inputs|}
\end{itemize}
\clearpage
\section*{GenAI interactions}
Below you will find all interactions between the `user' (program) and the `assistant' (OpenAI).
\subsection*{Program $\rightarrow$ GenAI}
\begin{lstlisting}
You are given the following task to perform in Dafny:

Write a method to satisfy the below requirements.

The signature should be:

p_4_2_c_cumulative_totals_seq(inputs: seq<int>) returns (totals: seq<int>)

The method should respect the following contract:

requires 0 < |inputs|, ensures |totals| == |inputs|, ensures forall i :: 0 <= i < |inputs| ==> totals[i] == (if i == 0 then inputs[0] else totals[i - 1] + inputs[i])

Produce and show only the Dafny body of this method, including the curly braces that surround it. Do not show the signature nor contract. You must not use 'assume' nor 'expect' anywhere in your code.
\end{lstlisting}
\subsection*{GenAI $\rightarrow$ Program}
\textbf{System fingerprint: }null
\\\textbf{ID: }chatcmpl-DnLKpSn6At7YODDa4zUcJ5Ytli0VN
\\\textbf{Model: }gpt-5.5-2026-04-23
\\\textbf{Created at: }1780652455
\\\textbf{Finish reason: }stop
\\\textbf{Usage: }358 tokens in, and 3124 tokens out
\begin{lstlisting}
{
  totals := [inputs[0]];
  var n := |inputs|;
  var k := 1;

  while k < n
    invariant n == |inputs|
    invariant 1 <= k <= n
    invariant |totals| == k
    invariant forall j :: 0 <= j < k ==> totals[j] == (if j == 0 then inputs[0] else totals[j - 1] + inputs[j])
    decreases n - k
  {
    var prevTotals := totals;
    assert |prevTotals| == k;
    assert forall j :: 0 <= j < k ==> prevTotals[j] == (if j == 0 then inputs[0] else prevTotals[j - 1] + inputs[j]);

    var next := prevTotals[k - 1] + inputs[k];
    totals := prevTotals + [next];

    forall j | 0 <= j < k
      ensures totals[j] == prevTotals[j]
    {
      assert 0 <= j < |prevTotals|;
      assert totals == prevTotals + [next];
      assert totals[j] == (prevTotals + [next])[j];
      assert (prevTotals + [next])[j] == prevTotals[j];
    }

    forall j | 0 <= j < k + 1
      ensures totals[j] == (if j == 0 then inputs[0] else totals[j - 1] + inputs[j])
    {
      if j < k {
        assert totals[j] == prevTotals[j];
        if j == 0 {
          assert prevTotals[j] == inputs[0];
        } else {
          assert 0 <= j - 1 < k;
          assert totals[j - 1] == prevTotals[j - 1];
          assert prevTotals[j] == prevTotals[j - 1] + inputs[j];
        }
      } else {
        assert j == k;
        assert 0 < j;
        assert j - 1 == k - 1;
        assert totals[j] == next;
        assert totals[j - 1] == prevTotals[k - 1];
        assert next == prevTotals[k - 1] + inputs[k];
      }
    }

    k := k + 1;
  }
}
\end{lstlisting}
\subsection*{Program $\rightarrow$ GenAI}
\begin{lstlisting}
When using dafny verify, the below error is emitted and verify fails:

problem_data/p_4_2_c_cumulative_totals_seq/program.dfy(39,10): Error: assertion might not hold
   |
39 |           assert prevTotals[j] == inputs[0];
   |           ^^^^^^


Dafny program verifier finished with 1 verified, 1 error

Correct the error in code by altering only the method body. Produce and show only the Dafny body, including the curly braces that surround it. Do not show the signature nor contract. You must not use 'assume' nor 'expect' anywhere in your code.
\end{lstlisting}
\subsection*{GenAI $\rightarrow$ Program}
\textbf{System fingerprint: }null
\\\textbf{ID: }chatcmpl-DnLM0tlHLFS8avdfY8Vn8bH6L6LHg
\\\textbf{Model: }gpt-5.5-2026-04-23
\\\textbf{Created at: }1780652528
\\\textbf{Finish reason: }stop
\\\textbf{Usage: }1016 tokens in, and 2873 tokens out
\begin{lstlisting}
{
  totals := [inputs[0]];
  var n := |inputs|;
  var k := 1;

  while k < n
    invariant n == |inputs|
    invariant 1 <= k <= n
    invariant |totals| == k
    invariant forall j :: 0 <= j < k ==> totals[j] == (if j == 0 then inputs[0] else totals[j - 1] + inputs[j])
    decreases n - k
  {
    var prevTotals := totals;
    assert |prevTotals| == k;
    assert forall j :: 0 <= j < k ==> prevTotals[j] == (if j == 0 then inputs[0] else prevTotals[j - 1] + inputs[j]);

    var next := prevTotals[k - 1] + inputs[k];
    totals := prevTotals + [next];

    forall j | 0 <= j < k
      ensures totals[j] == prevTotals[j]
    {
      assert 0 <= j < |prevTotals|;
      assert totals == prevTotals + [next];
      assert totals[j] == (prevTotals + [next])[j];
      assert (prevTotals + [next])[j] == prevTotals[j];
    }

    forall j | 0 <= j < k + 1
      ensures totals[j] == (if j == 0 then inputs[0] else totals[j - 1] + inputs[j])
    {
      if j < k {
        assert 0 <= j < k;
        assert totals[j] == prevTotals[j];
        assert prevTotals[j] == (if j == 0 then inputs[0] else prevTotals[j - 1] + inputs[j]);
        if j == 0 {
          assert prevTotals[j] == inputs[0];
        } else {
          assert 0 <= j - 1 < k;
          assert totals[j - 1] == prevTotals[j - 1];
          assert prevTotals[j] == prevTotals[j - 1] + inputs[j];
        }
      } else {
        assert j == k;
        assert 0 < j;
        assert j - 1 == k - 1;
        assert totals[j] == next;
        assert totals[j - 1] == prevTotals[k - 1];
        assert next == prevTotals[k - 1] + inputs[k];
      }
    }

    k := k + 1;
  }
}
\end{lstlisting}
\section*{Final Program}
\begin{lstlisting}
method p_4_2_c_cumulative_totals_seq(inputs: seq<int>) returns (totals: seq<int>)
	requires 0 < |inputs|
	ensures |totals| == |inputs|
	ensures forall i :: 0 <= i < |inputs| ==> totals[i] == (if i == 0 then inputs[0] else totals[i - 1] + inputs[i])
{
  totals := [inputs[0]];
  var n := |inputs|;
  var k := 1;

  while k < n
    invariant n == |inputs|
    invariant 1 <= k <= n
    invariant |totals| == k
    invariant forall j :: 0 <= j < k ==> totals[j] == (if j == 0 then inputs[0] else totals[j - 1] + inputs[j])
    decreases n - k
  {
    var prevTotals := totals;
    assert |prevTotals| == k;
    assert forall j :: 0 <= j < k ==> prevTotals[j] == (if j == 0 then inputs[0] else prevTotals[j - 1] + inputs[j]);

    var next := prevTotals[k - 1] + inputs[k];
    totals := prevTotals + [next];

    forall j | 0 <= j < k
      ensures totals[j] == prevTotals[j]
    {
      assert 0 <= j < |prevTotals|;
      assert totals == prevTotals + [next];
      assert totals[j] == (prevTotals + [next])[j];
      assert (prevTotals + [next])[j] == prevTotals[j];
    }

    forall j | 0 <= j < k + 1
      ensures totals[j] == (if j == 0 then inputs[0] else totals[j - 1] + inputs[j])
    {
      if j < k {
        assert 0 <= j < k;
        assert totals[j] == prevTotals[j];
        assert prevTotals[j] == (if j == 0 then inputs[0] else prevTotals[j - 1] + inputs[j]);
        if j == 0 {
          assert prevTotals[j] == inputs[0];
        } else {
          assert 0 <= j - 1 < k;
          assert totals[j - 1] == prevTotals[j - 1];
          assert prevTotals[j] == prevTotals[j - 1] + inputs[j];
        }
      } else {
        assert j == k;
        assert 0 < j;
        assert j - 1 == k - 1;
        assert totals[j] == next;
        assert totals[j - 1] == prevTotals[k - 1];
        assert next == prevTotals[k - 1] + inputs[k];
      }
    }

    k := k + 1;
  }
}
\end{lstlisting}
\section*{Total Token Usage}
\textbf{Input tokens: }1374
\\\textbf{Output tokens: }5997
\\\textbf{Reasoning tokens: }4807
\\\textbf{Sum of `total tokens': }7371
\section*{Experiment Timings}
\textbf{Overall Experiment} started at 1780652454921, ended at 1780652599541, lasting 144620ms (144.62 seconds)
\\\textbf{Iteration \#1} started at 1780652454921, ended at 1780652528040, lasting 73119ms (73.12 seconds)
\\\textbf{Iteration \#2} started at 1780652528040, ended at 1780652599540, lasting 71500ms (71.50 seconds)
\end{document}
